% Source: round3/exchanges/tensor_barrier.tex
\section{Tensor products preserve a three-row exchange barrier}
For a partial Hadamard matrix $P$, let $\rho(P)$ be the minimum number
of old rows deleted in an exchange which inserts one more row than it
deletes, all retained and inserted rows remaining mutually orthogonal.
Set $\rho(P)=\infty$ if there is no such exchange. The previous pass
certified $\rho(P)=3$ for one $16\times24$ matrix. Here we prove a
uniform preservation theorem, rather than search larger isolated orders.


\begin{proposition}[Tensoring cannot amplify the exchange cost upward]
If $H$ is a full Hadamard matrix, then
\[
 \rho(P\otimes H)\le\rho(P).
\]
\end{proposition}
\begin{proof}
Choose any row $h$ of $H$. A $k$-for-$(k+1)$ exchange of rows $p_i$
for rows $v_j$ lifts to an exchange of $p_i\otimes h$ for $v_j\otimes h$.
Rows in other Hadamard fibers are orthogonal to every inserted row,
while the same-fiber checks reduce to the original exchange.
\end{proof}
Thus repeated full-Hadamard tensoring alone cannot turn a finite
exchange lower bound into an unbounded one. The interesting issue is
whether the tensor permits a cheaper exchange by using non-product rows.

\subsection{A condition excluding cheaper non-product exchanges}
For a row-index set $I$, let $\mathcal C_I$ consist of the sign vectors
orthogonal to all rows of $P$ outside $I$. Sign vectors differing only
by a global minus sign are treated as one line. Write the order as $n$.

\begin{proposition}[A general span-rigidity criterion]
Suppose $P$ is maximal and, for every index set $I$ of size at most $c$,
\[
 \mathcal C_I\subseteq\operatorname{span}\{p_i:i\in I\}.
\]
Then no tensor $P\otimes H$ with a full Hadamard factor admits an
augmenting exchange deleting at most $c$ rows.
\end{proposition}
\begin{proof}
For a deletion set $D$ of size $k\le c$, let $I$ be its projection
onto the original row indices. Reshape a candidate as $X$. As in the
proof below, every column belongs to $\mathcal C_I$, so $X$ lies in
$\operatorname{span}(p_i:i\in I)\otimes\R^s$. Orthogonality to the
retained tensor rows then forces $X$ into the span of the $k$ deleted
tensor rows. That space has dimension $k$ and cannot contain $k+1$
mutually orthogonal nonzero rows.
\end{proof}
The order-twenty-four example fails this criterion at some pairs: its
two extra sign lines span a third direction. Their particular geometry,
rather than span rigidity, supplies the stronger preservation below.

\begin{theorem}[Preservation under half-overlap pair structure]
Suppose $P$ is maximal and has the following properties:
\begin{enumerate}
\item $\mathcal C_{\{i\}}$ has just the line of $p_i$.
\item For every pair $i,j$, either $\mathcal C_{\{i,j\}}$ has only the
lines of $p_i,p_j$, or it has exactly four lines. In the latter case
the two extra lines can be oriented as $x,y$ so that
\[
 x\cdot y=0,\qquad
 (p_i\cdot x,p_j\cdot x)=(p_i\cdot y,p_j\cdot y)
 =\tfrac n2(1,\sigma),\quad\sigma\in\{\pm1\}.
\]
\end{enumerate}
Then for every full Hadamard matrix $H$, $P\otimes H$ has no augmenting
exchange deleting at most two rows. In particular, if $\rho(P)=3$, then
$\rho(P\otimes H)=3$.
\end{theorem}
\begin{proof}
Let $H$ have order $s$, with rows $h_b$, and reshape a proposed sign
row of length $ns$ as an $n\times s$ sign matrix $X$. Its correlations
with the old tensor rows are the entries of $PXH^{\mathsf T}$.
If the deleted tensor rows project to original row-index set $I$,
every column of $X$ belongs to $\mathcal C_I$: for $k\notin I$,
$p_kXH^{\mathsf T}=0$ implies $p_kX=0$.

With no deletions, this contradicts maximality. With one deletion
$(i,b)$, every column is $\pm p_i$, so $X=p_i\otimes\varepsilon$;
the remaining fiber constraints force $\varepsilon$ to be proportional
to $h_b$. The only candidate line is the removed tensor row.

If both deleted rows have the same original index $i$ and different
fibers $b,c$, the same argument gives $X=p_i\otimes\varepsilon$, now
with $\varepsilon\in\operatorname{span}(h_b,h_c)\cap\{\pm1\}^s$.
That intersection has only $\pm h_b,\pm h_c$. Indeed write
$\varepsilon=\alpha h_b+\beta h_c$; the two orthogonal sign rows
agree and disagree on nonempty sets, so both $\alpha+\beta$ and
$\alpha-\beta$ are signs. Consequently one coefficient is zero and
the other is a sign. There are at most two candidate lines.

It remains to delete $(i,b),(j,c)$ with $i\ne j$. Fourier inversion
along the Hadamard factor gives
\[
 p_iX=\alpha h_b,\qquad p_jX=\beta h_c.
\]
For each column of $X$, the possible correlation pairs with $p_i,p_j$
are $\pm(n,0)$, $\pm(0,n)$ and, when present,
$\pm(n/2)(1,\sigma)$. Their absolute magnitudes are constant across
columns, because $h_b,h_c$ have sign entries. If one of $\alpha,\beta$
vanishes, the candidate is one of the removed old tensor rows. If both
are nonzero, the extra pair type is necessary, and its fixed correlation
ratio forces $h_b$ and $h_c$ to be proportional. Thus $b=c$; otherwise
there are again only the two removed candidate lines.

In the remaining case put $h=h_b=h_c$ and orient globally so that the
extra correlation pair is $(n/2)(1,\sigma)h$.
Let
\[
 w=(p_i+\sigma p_j)/2,\qquad z=(x-y)/2.
\]
Projection onto $\operatorname{span}(p_i,p_j)$ shows that
$x=w+z$, $y=w-z$, $w\perp z$, and
$\|w\|^2=\|z\|^2=n/2$. To justify the residual cancellation, the
residuals of $x,y$ each have squared norm $n/2$ and inner product
$-n/2$, so they are negatives of each other. Every extra tensor
candidate therefore has the form
\[
 W_{\varepsilon}=w\otimes h+z\otimes\varepsilon,
 \qquad\varepsilon\in\{\pm1\}^s.
\]
Neither removed old tensor row is orthogonal to any $W_\varepsilon$.
For two extra candidates,
\[
 W_\varepsilon\cdot W_\delta
 =\tfrac n2(s+\varepsilon\cdot\delta),
\]
which is zero exactly when $\delta=-\varepsilon$. Thus the candidate
orthogonality graph consists of the old orthogonal pair and a matching
on the extra candidates. It has no triangle, so no three mutually
orthogonal replacement rows exist. All cases are exhausted.
The upper bound from the preceding proposition proves equality when
$\rho(P)=3$.
\end{proof}

\subsection{An infinite family with exact exchange cost three}
The independently certified $16\times24$ matrix from the previous pass
satisfies the hypotheses. Among its $120$ two-row deletions, one hundred
have only the two old candidate lines. The other twenty have exactly
two extra lines. Those extras are mutually orthogonal, their
correlations with the two old rows are all $\pm12$, and the two
correlation vectors are parallel. All these identities are checked
exactly from the exhaustive sign-cube certificate.

\begin{corollary}
For every Hadamard order $s$, there is a maximal partial matrix of
shape $16s\times24s$ whose minimum augmenting exchange deletes exactly
three rows. In particular such examples exist at every order
$24\cdot2^a$.
\end{corollary}

This is a statement about partial matrices and the cost of escaping
their terminal configurations. It does not produce new Hadamard orders.
The base three-for-four exchange has an explicit full completion. Apply
that exchange separately in each of the $s$ Hadamard fibers. After
$s$ exchanges the result is $Q\otimes H$, where the base $17$-row
matrix $Q$ completes; hence this whole tensor family can be completed
after those $s$ exchanges. A single lifted exchange is only claimed
to augment the tensor partial matrix, not necessarily complete it.

\paragraph{Exact checks.}
The script \path{round3/exchanges/check.py} re-enumerates all $2^{23}$
normalized sign vectors for the base matrix, verifies the pair-structure
hypotheses for all $120$ pairs, and checks lifted exchanges and
completions at tensor-factor orders $1,2,4,8$. Additional exact
low-dimensional nullspace enumeration at factor order two confirms
the predicted candidate sets in each of the three deletion cases.
The uniform statement rests on the proof above. This statement is not claimed as Lean verified.

